\subsection{Paths as Structured Filesystem References}

A path is not the file itself. A path is a reference that tells the operating system where a filesystem object should be found or created. That object may be a file, a directory, a symbolic link, or some other platform-specific filesystem entry.

Earlier file examples used simple string paths such as \texttt{"demo.txt"}. That works, but a path has internal structure: parent directories, final names, suffixes, roots, drives, separators, and resolution rules. Python can represent paths as ordinary strings, but it also provides \texttt{pathlib.Path}, which exposes that structure through object methods and attributes.

\subsubsection{String Paths vs. \texttt{pathlib.Path} Objects}

A string path is just a \texttt{str} object whose content happens to describe a filesystem location. Python does not automatically understand its pieces as path components. It can still be passed to \texttt{open()}, but string methods such as \texttt{split()} and \texttt{replace()} are general text methods, not filesystem-aware path operations.

\begin{verbatim}
from pathlib import Path

str_path = "data/chuck_facts.txt"
path_obj = Path("data") / "chuck_facts.txt"

print("Two path representations:")
print(f"str_path = {str_path}")
print(f"path_obj = {path_obj}")

print()
print("Object types:")
print(f"type(str_path) = {type(str_path)}")
print(f"type(path_obj) = {type(path_obj)}")

print()
print("Selected names from dir(str_path):")
for name in ["split", "replace", "startswith", "endswith", "upper"]:
    print(f"{name:12} -> {name in dir(str_path)}")

print()
print("Selected names from dir(path_obj):")
for name in ["name", "parent", "suffix", "stem", "parts", "exists", "is_file"]:
    print(f"{name:12} -> {name in dir(path_obj)}")

print()
print("Path object components:")
print(f"path_obj.name   = {path_obj.name}")
print(f"path_obj.parent = {path_obj.parent}")
print(f"path_obj.stem   = {path_obj.stem}")
print(f"path_obj.suffix = {path_obj.suffix}")
print(f"path_obj.parts  = {path_obj.parts}")
\end{verbatim}

The string path and the \texttt{Path} object may print similarly, but they are not the same kind of object. The string exposes text behavior. The \texttt{Path} object exposes filesystem-reference behavior.

The operator \texttt{/} is overloaded for \texttt{Path} objects. It does not mean division here. It means path composition. The expression \texttt{Path("data") / "chuck\_facts.txt"} creates a new path by joining a directory-like component with a final filename component.

A \texttt{Path} object can still be converted to a string when necessary.

\begin{verbatim}
from pathlib import Path

path_obj = Path("quotes") / "monty_python.txt"
path_text = str(path_obj)

print("Path object:")
print(f"path_obj       = {path_obj}")
print(f"type(path_obj) = {type(path_obj)}")

print()
print("Converted string:")
print(f"path_text       = {path_text}")
print(f"type(path_text) = {type(path_text)}")
\end{verbatim}

For modern Python code, \texttt{Path} objects are usually clearer than manually assembled string paths. They keep path construction separate from ordinary text manipulation.

\subsubsection{Absolute Paths, Relative Paths, and Current Working Directory Resolution}

A relative path is interpreted from some starting directory. For ordinary script execution, that starting point is usually the current working directory. An absolute path already contains enough information to identify a location from the filesystem root or drive root.

This distinction matters because the same relative path can point to different files when the program is started from different directories.

\begin{verbatim}
from pathlib import Path

cwd = Path.cwd()

rel_path = Path("data") / "answer.txt"
abs_path = rel_path.resolve()

print("Current working directory:")
print(f"cwd       = {cwd}")
print(f"type(cwd) = {type(cwd)}")

print()
print("Relative path:")
print(f"rel_path       = {rel_path}")
print(f"type(rel_path) = {type(rel_path)}")
print(f"rel_path.is_absolute() = {rel_path.is_absolute()}")

print()
print("Resolved absolute path:")
print(f"abs_path       = {abs_path}")
print(f"type(abs_path) = {type(abs_path)}")
print(f"abs_path.is_absolute() = {abs_path.is_absolute()}")

print()
print("Selected names from dir(rel_path):")
for name in ["resolve", "absolute", "is_absolute", "parent", "name", "parts"]:
    print(f"{name:14} -> {name in dir(rel_path)}")
\end{verbatim}

The method \texttt{Path.cwd()} returns the current working directory as a \texttt{Path} object. The method \texttt{resolve()} returns an absolute version of a path. In normal use, this means the relative path is interpreted against the current working directory.

A relative path is not incomplete text. It is a valid reference whose meaning depends on where the process is currently positioned in the filesystem.

\begin{verbatim}
from pathlib import Path

rel_path = Path("logs") / "seinfeld.txt"

print("Relative path parts:")
print(f"rel_path.parts = {rel_path.parts}")

print()
print("If the current working directory changes,")
print("the same relative path may refer to another external location.")

print()
print(f"current directory = {Path.cwd()}")
print(f"relative path     = {rel_path}")
print(f"resolved path     = {rel_path.resolve()}")
\end{verbatim}

The practical rule is simple: relative paths are convenient inside a known project directory; absolute paths are clearer when the program must identify one specific location regardless of where it was launched.

\subsubsection{Platform Separators: Windows Backslashes, POSIX Slashes, and Portable Path Construction}

Different operating systems display paths differently. POSIX-style systems, such as Linux and macOS, normally use forward slashes. Windows paths commonly use backslashes and may include drive names.

Hard-coding separators manually is fragile. Python can build portable paths with \texttt{pathlib.Path} instead.

\begin{verbatim}
from pathlib import Path, PureWindowsPath, PurePosixPath

portable_path = Path("course_data") / "quotes" / "simpsons.txt"

print("Portable path built with pathlib:")
print(f"portable_path       = {portable_path}")
print(f"type(portable_path) = {type(portable_path)}")
print(f"portable_path.parts = {portable_path.parts}")

print()
print("Pure path examples, without touching the real filesystem:")

win_path = PureWindowsPath("course_data") / "quotes" / "simpsons.txt"
posix_path = PurePosixPath("course_data") / "quotes" / "simpsons.txt"

print(f"win_path       = {win_path}")
print(f"type(win_path) = {type(win_path)}")
print(f"win_path.parts = {win_path.parts}")

print()
print(f"posix_path       = {posix_path}")
print(f"type(posix_path) = {type(posix_path)}")
print(f"posix_path.parts = {posix_path.parts}")
\end{verbatim}

\texttt{PureWindowsPath} and \texttt{PurePosixPath} are useful for demonstrating path syntax rules without actually requiring the program to run on that operating system. They model path structure, but they do not access the filesystem.

Backslashes inside ordinary Python string literals also have a second problem: backslash begins escape sequences such as \texttt{\textbackslash n} and \texttt{\textbackslash t}. This can accidentally corrupt a Windows-style path if it is written carelessly.

\begin{verbatim}
bad_path = "C:\new\text.txt"
raw_path = r"C:\new\text.txt"

print("Careless Windows-style string:")
print(bad_path)
print(repr(bad_path))

print()
print("Raw string version:")
print(raw_path)
print(repr(raw_path))
\end{verbatim}

In the first string, \texttt{\textbackslash n} becomes a newline and \texttt{\textbackslash t} becomes a tab. The string no longer contains the visible path the programmer probably intended. A raw string avoids that escape processing, but \texttt{pathlib} is usually cleaner because it reduces manual separator work.

\begin{verbatim}
from pathlib import Path

path_obj = Path("C:/new") / "text.txt"

print("Path constructed without manual backslash escaping:")
print(path_obj)
print(f"type(path_obj) = {type(path_obj)}")
\end{verbatim}

The practical rule is direct: use \texttt{Path(...)} and the \texttt{/} composition operator for path construction. Let Python represent the path according to the current platform, and avoid treating filesystem references as ordinary string concatenation problems.